%
% claude-code-lux.tex
%
% Z Specification: Claude Code + Lux Visual Surface State (Layer 2)
% Extends the base Claude Code state machine with lux display state and
% consumer integration patterns for beads, quarry, biff, and z-spec.
%
% Type-check: fuzz -t claude-code-lux.tex
% Animate:    probcli claude-code-lux.tex -animate 20
%

\documentclass[a4paper,10pt,fleqn]{article}
\usepackage[margin=1in]{geometry}
\usepackage{fuzz}
\usepackage[colorlinks=true,linkcolor=blue,citecolor=blue,urlcolor=blue]{hyperref}

\begin{document}

\title{Claude Code + Lux: A Z Specification}
\author{Layer 2 --- Visual Output Surface State}
\date{March 2026}
\hypersetup{
  pdftitle={Claude Code + Lux: A Z Specification},
  pdfauthor={Punt Labs},
  pdfsubject={Z Specification},
  pdfcreator={fuzz/probcli}
}
\maketitle

\tableofcontents
\newpage

% ============================================================
\section{Introduction}
% ============================================================

This specification extends the base Claude Code state machine
(\texttt{claude-code.tex}) with lux visual output surface state.
Lux is a \emph{building block} (see \texttt{integration.md}): it
provides deterministic rendering at L1 that agents compose at L4.
Lux has no upstream awareness --- it does not know about beads,
biff, quarry, or any consumer. Consumers discover lux at L1 and
use it to render visual output.

The model captures two concerns:

\begin{enumerate}
  \item \textbf{Lux display state} --- whether the display is active,
    what scene is shown, whether an interaction is pending.
  \item \textbf{Consumer integration} --- how other tools trigger
    lux displays via hooks and skills, gated by a display mode flag.
\end{enumerate}

The key properties verified:

\begin{enumerate}
  \item \textbf{Scene requires display}: a scene can only be active
    when the display is enabled.
  \item \textbf{No orphan interactions}: blocking for user interaction
    (\texttt{recv}) requires an active scene.
  \item \textbf{Display mode gates rendering}: when
    \texttt{displayMode = dmOff}, no consumer triggers lux displays.
  \item \textbf{Scene updates are safe}: \texttt{UpdateScene} requires
    an active scene. No updates to a cleared display.
  \item \textbf{Beads board refresh}: after bead state changes
    (\texttt{ClaimBead}, \texttt{CloseBead}), the beads board is
    marked stale for refresh.
\end{enumerate}

% ============================================================
\section{Base Types (from claude-code.tex)}
% ============================================================

\begin{zed}
  [TOOLNAME, TOOLINPUT, TOOLOUTPUT, PROMPT, RESPONSE, AGENTID, SESSIONID, CONTEXTCHUNK]
\end{zed}

\begin{zed}
  SessionPhase ::= spInactive | spIdle | spProcessing | spResponding | spEnded
\end{zed}

\begin{zed}
  ToolPhase ::= tpNone | tpPreHook | tpPermission | tpExecuting | tpPostHook
\end{zed}

\begin{zed}
  PermissionMode ::= pmDefault | pmPlan | pmAcceptEdits | pmDontAsk | pmBypass
\end{zed}

\begin{zed}
  PermissionDecision ::= pdAllow | pdDeny | pdAsk
\end{zed}

\begin{zed}
  HookDecision ::= hdAllow | hdBlock
\end{zed}

\begin{zed}
  SessionStartSource ::= ssStartup | ssResume | ssClear | ssCompact
\end{zed}

\begin{zed}
  SessionEndReason ::= seClear | seLogout | sePromptExit | seBypassDisabled | seOther
\end{zed}

\begin{zed}
  ZBOOL ::= ztrue | zfalse
\end{zed}

% ============================================================
\section{Lux Types}
% ============================================================

\begin{zed}
  [SCENEID, ELEMENTID]
\end{zed}

\texttt{SCENEID} identifies a displayed scene (e.g.,
\texttt{beads-board}, \texttt{quarry-search}, \texttt{z-spec-states}).
\texttt{ELEMENTID} identifies an interactive element within a scene
(e.g., a button, slider, or table row).

\begin{zed}
  DisplayMode ::= dmOff | dmOn
\end{zed}

Whether lux visual output is enabled. \texttt{dmOn} means consumers
should render visual output when appropriate. \texttt{dmOff} means
consumers degrade to text-only. Analogous to vox's
\texttt{notifyMode}.

\begin{zed}
  SceneKind ::= skBeads | skDashboard | skDataExplorer | skDiagram | skCustom
\end{zed}

Classification of the active scene. Used by consumer hooks to decide
whether to update or replace the display.

% ============================================================
\section{Global Constants}
% ============================================================

\begin{axdef}
  maxTools : \nat \\
  maxSubagents : \nat \\
  maxContext : \nat
  \where
  maxTools = 5 \\
  maxSubagents = 2 \\
  maxContext = 3
\end{axdef}

% ============================================================
\section{State}
% ============================================================

\begin{schema}{State}
  %
  % --- Base Claude Code state ---
  %
  sessionId : \finset SESSIONID \\
  sessionPhase : SessionPhase \\
  permMode : PermissionMode \\
  toolPhase : ToolPhase \\
  currentTool : \finset TOOLNAME \\
  toolsThisTurn : \nat \\
  activeSubagents : \finset AGENTID \\
  context : \seq CONTEXTCHUNK \\
  pendingPrompt : \finset PROMPT \\
  pendingResponse : \finset RESPONSE \\
  %
  % --- Lux display state ---
  %
  displayMode : DisplayMode \\
  activeScene : \finset SCENEID \\
  sceneKind : \finset SceneKind \\
  pendingRecv : ZBOOL \\
  beadsBoardStale : ZBOOL
  \where
  %
  % === Base invariants ===
  %
  \# sessionId \leq 1 \\
  sessionPhase = spInactive \implies sessionId = \emptyset \\
  sessionPhase \neq spInactive \implies \# sessionId = 1 \\
  sessionPhase \neq spProcessing \implies toolPhase = tpNone \\
  \# currentTool \leq 1 \\
  toolPhase = tpNone \implies currentTool = \emptyset \\
  toolPhase \neq tpNone \implies \# currentTool = 1 \\
  toolsThisTurn \leq maxTools \\
  sessionPhase = spInactive \implies activeSubagents = \emptyset \\
  sessionPhase = spEnded \implies activeSubagents = \emptyset \\
  \# activeSubagents \leq maxSubagents \\
  \# context \leq maxContext \\
  \# pendingPrompt \leq 1 \\
  sessionPhase = spIdle \implies pendingPrompt = \emptyset \\
  sessionPhase = spInactive \implies pendingPrompt = \emptyset \\
  sessionPhase = spEnded \implies pendingPrompt = \emptyset \\
  \# pendingResponse \leq 1 \\
  sessionPhase \neq spResponding \implies pendingResponse = \emptyset \\
  %
  % === Lux invariants ===
  %
  % At most one active scene
  %
  \# activeScene \leq 1 \\
  \# sceneKind \leq 1 \\
  %
  % Scene requires display enabled
  %
  displayMode = dmOff \implies activeScene = \emptyset \\
  %
  % Scene kind tracks active scene
  %
  activeScene = \emptyset \implies sceneKind = \emptyset \\
  activeScene \neq \emptyset \implies \# sceneKind = 1 \\
  %
  % Pending interaction requires active scene
  %
  pendingRecv = ztrue \implies activeScene \neq \emptyset \\
  %
  % Beads board staleness only when beads board is active
  %
  beadsBoardStale = ztrue \implies sceneKind = \{ skBeads \} \\
  %
  % No lux state when session is inactive
  %
  sessionPhase = spInactive \implies displayMode = dmOff \\
  sessionPhase = spInactive \implies pendingRecv = zfalse \\
  sessionPhase = spInactive \implies beadsBoardStale = zfalse
\end{schema}

% ============================================================
\section{Initialization}
% ============================================================

\begin{schema}{Init}
  State'
  \where
  sessionId' = \emptyset \\
  sessionPhase' = spInactive \\
  permMode' = pmDefault \\
  toolPhase' = tpNone \\
  currentTool' = \emptyset \\
  toolsThisTurn' = 0 \\
  activeSubagents' = \emptyset \\
  context' = \langle \rangle \\
  pendingPrompt' = \emptyset \\
  pendingResponse' = \emptyset \\
  %
  displayMode' = dmOff \\
  activeScene' = \emptyset \\
  sceneKind' = \emptyset \\
  pendingRecv' = zfalse \\
  beadsBoardStale' = zfalse
\end{schema}

% ============================================================
\section{Lux Frame Condition}
% ============================================================

\begin{schema}{LuxFrame}
  \Delta State
  \where
  displayMode' = displayMode \\
  activeScene' = activeScene \\
  sceneKind' = sceneKind \\
  pendingRecv' = pendingRecv \\
  beadsBoardStale' = beadsBoardStale
\end{schema}

% ============================================================
\section{Session Lifecycle Operations}
% ============================================================

\subsection{StartSession}

Lux reads display mode from config (\texttt{.lux/config.md} or
equivalent sentinel). If no config exists, defaults to \texttt{dmOff}.

\begin{schema}{StartSession}
  \Delta State \\
  sid? : SESSIONID \\
  source? : SessionStartSource \\
  mode? : PermissionMode \\
  initDisplay? : DisplayMode
  \where
  sessionPhase = spInactive \\
  sessionPhase' = spIdle \\
  sessionId' = \{ sid? \} \\
  permMode' = mode? \\
  toolPhase' = tpNone \\
  currentTool' = \emptyset \\
  toolsThisTurn' = 0 \\
  activeSubagents' = \emptyset \\
  context' = \langle \rangle \\
  pendingPrompt' = \emptyset \\
  pendingResponse' = \emptyset \\
  %
  displayMode' = initDisplay? \\
  activeScene' = \emptyset \\
  sceneKind' = \emptyset \\
  pendingRecv' = zfalse \\
  beadsBoardStale' = zfalse
\end{schema}

\subsection{SubmitPrompt}

\begin{schema}{SubmitPrompt}
  LuxFrame \\
  prompt? : PROMPT \\
  hookResult? : HookDecision \\
  contextChunk? : CONTEXTCHUNK
  \where
  sessionPhase = spIdle \\
  hookResult? = hdAllow \\
  sessionPhase' = spProcessing \\
  pendingPrompt' = \{ prompt? \} \\
  context' = (1 \upto maxContext) \dres (context \cat \langle contextChunk? \rangle) \\
  toolsThisTurn' = 0 \\
  toolPhase' = tpNone \\
  currentTool' = \emptyset \\
  sessionId' = sessionId \\
  permMode' = permMode \\
  activeSubagents' = activeSubagents \\
  pendingResponse' = \emptyset
\end{schema}

\begin{schema}{RejectPrompt}
  \Xi State \\
  prompt? : PROMPT \\
  hookResult? : HookDecision
  \where
  sessionPhase = spIdle \\
  hookResult? = hdBlock
\end{schema}

% ============================================================
\section{Tool Call Operations}
% ============================================================

\subsection{BeginToolCall}

\begin{schema}{BeginToolCall}
  LuxFrame \\
  tool? : TOOLNAME
  \where
  sessionPhase = spProcessing \\
  toolPhase = tpNone \\
  toolsThisTurn < maxTools \\
  toolPhase' = tpPreHook \\
  currentTool' = \{ tool? \} \\
  sessionPhase' = sessionPhase \\
  sessionId' = sessionId \\
  permMode' = permMode \\
  toolsThisTurn' = toolsThisTurn \\
  activeSubagents' = activeSubagents \\
  context' = context \\
  pendingPrompt' = pendingPrompt \\
  pendingResponse' = pendingResponse
\end{schema}

\subsection{PreToolHookAllow}

\begin{schema}{PreToolHookAllow}
  LuxFrame \\
  hookDecision? : PermissionDecision \\
  needsPermission? : ZBOOL
  \where
  sessionPhase = spProcessing \\
  toolPhase = tpPreHook \\
  hookDecision? = pdAllow \\
  toolPhase' = \IF needsPermission? = ztrue \THEN tpPermission \ELSE tpExecuting \\
  sessionPhase' = sessionPhase \\
  sessionId' = sessionId \\
  permMode' = permMode \\
  currentTool' = currentTool \\
  toolsThisTurn' = toolsThisTurn \\
  activeSubagents' = activeSubagents \\
  context' = context \\
  pendingPrompt' = pendingPrompt \\
  pendingResponse' = pendingResponse
\end{schema}

\subsection{PreToolHookDeny}

\begin{schema}{PreToolHookDeny}
  LuxFrame \\
  hookDecision? : PermissionDecision \\
  denialContext? : CONTEXTCHUNK
  \where
  sessionPhase = spProcessing \\
  toolPhase = tpPreHook \\
  hookDecision? = pdDeny \\
  toolPhase' = tpNone \\
  currentTool' = \emptyset \\
  context' = (1 \upto maxContext) \dres (context \cat \langle denialContext? \rangle) \\
  sessionPhase' = sessionPhase \\
  sessionId' = sessionId \\
  permMode' = permMode \\
  toolsThisTurn' = toolsThisTurn \\
  activeSubagents' = activeSubagents \\
  pendingPrompt' = pendingPrompt \\
  pendingResponse' = pendingResponse
\end{schema}

\subsection{GrantPermission}

\begin{schema}{GrantPermission}
  LuxFrame
  \where
  sessionPhase = spProcessing \\
  toolPhase = tpPermission \\
  toolPhase' = tpExecuting \\
  sessionPhase' = sessionPhase \\
  sessionId' = sessionId \\
  permMode' = permMode \\
  currentTool' = currentTool \\
  toolsThisTurn' = toolsThisTurn \\
  activeSubagents' = activeSubagents \\
  context' = context \\
  pendingPrompt' = pendingPrompt \\
  pendingResponse' = pendingResponse
\end{schema}

\subsection{DenyPermission}

\begin{schema}{DenyPermission}
  LuxFrame \\
  denialContext? : CONTEXTCHUNK
  \where
  sessionPhase = spProcessing \\
  toolPhase = tpPermission \\
  toolPhase' = tpNone \\
  currentTool' = \emptyset \\
  context' = (1 \upto maxContext) \dres (context \cat \langle denialContext? \rangle) \\
  sessionPhase' = sessionPhase \\
  sessionId' = sessionId \\
  permMode' = permMode \\
  toolsThisTurn' = toolsThisTurn \\
  activeSubagents' = activeSubagents \\
  pendingPrompt' = pendingPrompt \\
  pendingResponse' = pendingResponse
\end{schema}

\subsection{CompleteToolSuccess}

\begin{schema}{CompleteToolSuccess}
  LuxFrame \\
  output? : CONTEXTCHUNK
  \where
  sessionPhase = spProcessing \\
  toolPhase = tpExecuting \\
  toolPhase' = tpNone \\
  currentTool' = \emptyset \\
  toolsThisTurn' = toolsThisTurn + 1 \\
  context' = (1 \upto maxContext) \dres (context \cat \langle output? \rangle) \\
  sessionPhase' = sessionPhase \\
  sessionId' = sessionId \\
  permMode' = permMode \\
  activeSubagents' = activeSubagents \\
  pendingPrompt' = pendingPrompt \\
  pendingResponse' = pendingResponse
\end{schema}

\subsection{CompleteToolFailure}

\begin{schema}{CompleteToolFailure}
  LuxFrame \\
  errorContext? : CONTEXTCHUNK
  \where
  sessionPhase = spProcessing \\
  toolPhase = tpExecuting \\
  toolPhase' = tpNone \\
  currentTool' = \emptyset \\
  toolsThisTurn' = toolsThisTurn + 1 \\
  context' = (1 \upto maxContext) \dres (context \cat \langle errorContext? \rangle) \\
  sessionPhase' = sessionPhase \\
  sessionId' = sessionId \\
  permMode' = permMode \\
  activeSubagents' = activeSubagents \\
  pendingPrompt' = pendingPrompt \\
  pendingResponse' = pendingResponse
\end{schema}

% ============================================================
\section{Lux Display Operations}
% ============================================================

\subsection{EnableDisplay}

User calls \texttt{/lux y}. Enables visual output mode.

\begin{schema}{EnableDisplay}
  \Delta State
  \where
  sessionPhase = spProcessing \\
  toolPhase = tpNone \\
  %
  displayMode' = dmOn \\
  %
  sessionPhase' = sessionPhase \\
  sessionId' = sessionId \\
  permMode' = permMode \\
  toolPhase' = toolPhase \\
  currentTool' = currentTool \\
  toolsThisTurn' = toolsThisTurn \\
  activeSubagents' = activeSubagents \\
  context' = context \\
  pendingPrompt' = pendingPrompt \\
  pendingResponse' = pendingResponse \\
  activeScene' = activeScene \\
  sceneKind' = sceneKind \\
  pendingRecv' = pendingRecv \\
  beadsBoardStale' = beadsBoardStale
\end{schema}

\subsection{DisableDisplay}

User calls \texttt{/lux n}. Disables visual output and clears
the display.

\begin{schema}{DisableDisplay}
  \Delta State
  \where
  sessionPhase = spProcessing \\
  toolPhase = tpNone \\
  %
  displayMode' = dmOff \\
  activeScene' = \emptyset \\
  sceneKind' = \emptyset \\
  pendingRecv' = zfalse \\
  beadsBoardStale' = zfalse \\
  %
  sessionPhase' = sessionPhase \\
  sessionId' = sessionId \\
  permMode' = permMode \\
  toolPhase' = toolPhase \\
  currentTool' = currentTool \\
  toolsThisTurn' = toolsThisTurn \\
  activeSubagents' = activeSubagents \\
  context' = context \\
  pendingPrompt' = pendingPrompt \\
  pendingResponse' = pendingResponse
\end{schema}

\subsection{ShowScene}

Claude calls \texttt{lux show()} to display a scene. Requires
display mode enabled.

\begin{schema}{ShowScene}
  \Delta State \\
  scene? : SCENEID \\
  kind? : SceneKind
  \where
  sessionPhase = spProcessing \\
  toolPhase = tpNone \\
  displayMode = dmOn \\
  %
  activeScene' = \{ scene? \} \\
  sceneKind' = \{ kind? \} \\
  pendingRecv' = zfalse \\
  beadsBoardStale' = \IF kind? = skBeads \THEN zfalse \ELSE beadsBoardStale \\
  %
  sessionPhase' = sessionPhase \\
  sessionId' = sessionId \\
  permMode' = permMode \\
  toolPhase' = toolPhase \\
  currentTool' = currentTool \\
  toolsThisTurn' = toolsThisTurn \\
  activeSubagents' = activeSubagents \\
  context' = context \\
  pendingPrompt' = pendingPrompt \\
  pendingResponse' = pendingResponse \\
  displayMode' = displayMode
\end{schema}

\subsection{UpdateScene}

Claude calls \texttt{lux update()} to patch the active scene.
Requires an active scene.

\begin{schema}{UpdateScene}
  \Delta State
  \where
  sessionPhase = spProcessing \\
  toolPhase = tpNone \\
  activeScene \neq \emptyset \\
  %
  % Scene remains active; staleness may be cleared if beads refresh
  beadsBoardStale' = zfalse \\
  %
  sessionPhase' = sessionPhase \\
  sessionId' = sessionId \\
  permMode' = permMode \\
  toolPhase' = toolPhase \\
  currentTool' = currentTool \\
  toolsThisTurn' = toolsThisTurn \\
  activeSubagents' = activeSubagents \\
  context' = context \\
  pendingPrompt' = pendingPrompt \\
  pendingResponse' = pendingResponse \\
  displayMode' = displayMode \\
  activeScene' = activeScene \\
  sceneKind' = sceneKind \\
  pendingRecv' = pendingRecv
\end{schema}

\subsection{ClearDisplay}

Claude calls \texttt{lux clear()}.

\begin{schema}{ClearDisplay}
  \Delta State
  \where
  sessionPhase = spProcessing \\
  toolPhase = tpNone \\
  %
  activeScene' = \emptyset \\
  sceneKind' = \emptyset \\
  pendingRecv' = zfalse \\
  beadsBoardStale' = zfalse \\
  %
  sessionPhase' = sessionPhase \\
  sessionId' = sessionId \\
  permMode' = permMode \\
  toolPhase' = toolPhase \\
  currentTool' = currentTool \\
  toolsThisTurn' = toolsThisTurn \\
  activeSubagents' = activeSubagents \\
  context' = context \\
  pendingPrompt' = pendingPrompt \\
  pendingResponse' = pendingResponse \\
  displayMode' = displayMode
\end{schema}

\subsection{WaitForInteraction}

Claude calls \texttt{lux recv()} to block for user interaction.

\begin{schema}{WaitForInteraction}
  \Delta State
  \where
  sessionPhase = spProcessing \\
  toolPhase = tpNone \\
  activeScene \neq \emptyset \\
  pendingRecv = zfalse \\
  %
  pendingRecv' = ztrue \\
  %
  sessionPhase' = sessionPhase \\
  sessionId' = sessionId \\
  permMode' = permMode \\
  toolPhase' = toolPhase \\
  currentTool' = currentTool \\
  toolsThisTurn' = toolsThisTurn \\
  activeSubagents' = activeSubagents \\
  context' = context \\
  pendingPrompt' = pendingPrompt \\
  pendingResponse' = pendingResponse \\
  displayMode' = displayMode \\
  activeScene' = activeScene \\
  sceneKind' = sceneKind \\
  beadsBoardStale' = beadsBoardStale
\end{schema}

\subsection{ReceiveInteraction}

User interacts with the display (button click, row selection, etc.).

\begin{schema}{ReceiveInteraction}
  \Delta State \\
  interactionContext? : CONTEXTCHUNK
  \where
  sessionPhase = spProcessing \\
  pendingRecv = ztrue \\
  %
  pendingRecv' = zfalse \\
  context' = (1 \upto maxContext) \dres (context \cat \langle interactionContext? \rangle) \\
  %
  sessionPhase' = sessionPhase \\
  sessionId' = sessionId \\
  permMode' = permMode \\
  toolPhase' = toolPhase \\
  currentTool' = currentTool \\
  toolsThisTurn' = toolsThisTurn \\
  activeSubagents' = activeSubagents \\
  pendingPrompt' = pendingPrompt \\
  pendingResponse' = pendingResponse \\
  displayMode' = displayMode \\
  activeScene' = activeScene \\
  sceneKind' = sceneKind \\
  beadsBoardStale' = beadsBoardStale
\end{schema}

% ============================================================
\section{Consumer Integration Operations}
% ============================================================

These operations model how other tools trigger lux displays.
They are gated by \texttt{displayMode = dmOn} --- when lux is
disabled, consumers degrade to text-only.

\subsection{MarkBeadsBoardStale}

A bead state change occurs (\texttt{bd update}, \texttt{bd close},
\texttt{bd create}). If the beads board is the active scene,
mark it stale so Claude knows to refresh.

\begin{schema}{MarkBeadsBoardStale}
  \Delta State
  \where
  sessionPhase = spProcessing \\
  toolPhase = tpNone \\
  displayMode = dmOn \\
  sceneKind = \{ skBeads \} \\
  %
  beadsBoardStale' = ztrue \\
  %
  sessionPhase' = sessionPhase \\
  sessionId' = sessionId \\
  permMode' = permMode \\
  toolPhase' = toolPhase \\
  currentTool' = currentTool \\
  toolsThisTurn' = toolsThisTurn \\
  activeSubagents' = activeSubagents \\
  context' = context \\
  pendingPrompt' = pendingPrompt \\
  pendingResponse' = pendingResponse \\
  displayMode' = displayMode \\
  activeScene' = activeScene \\
  sceneKind' = sceneKind \\
  pendingRecv' = pendingRecv
\end{schema}

% ============================================================
\section{Subagent Operations}
% ============================================================

\begin{schema}{SpawnSubagent}
  LuxFrame \\
  agentId? : AGENTID
  \where
  sessionPhase = spProcessing \\
  agentId? \notin activeSubagents \\
  \# activeSubagents < maxSubagents \\
  activeSubagents' = activeSubagents \cup \{ agentId? \} \\
  sessionPhase' = sessionPhase \\
  sessionId' = sessionId \\
  permMode' = permMode \\
  toolPhase' = toolPhase \\
  currentTool' = currentTool \\
  toolsThisTurn' = toolsThisTurn \\
  context' = context \\
  pendingPrompt' = pendingPrompt \\
  pendingResponse' = pendingResponse
\end{schema}

\begin{schema}{StopSubagent}
  LuxFrame \\
  agentId? : AGENTID \\
  hookResult? : HookDecision \\
  resultContext? : CONTEXTCHUNK
  \where
  sessionPhase = spProcessing \\
  agentId? \in activeSubagents \\
  hookResult? = hdAllow \\
  activeSubagents' = activeSubagents \setminus \{ agentId? \} \\
  context' = (1 \upto maxContext) \dres (context \cat \langle resultContext? \rangle) \\
  sessionPhase' = sessionPhase \\
  sessionId' = sessionId \\
  permMode' = permMode \\
  toolPhase' = toolPhase \\
  currentTool' = currentTool \\
  toolsThisTurn' = toolsThisTurn \\
  pendingPrompt' = pendingPrompt \\
  pendingResponse' = pendingResponse
\end{schema}

% ============================================================
\section{Response, Stop, Context, and Session End}
% ============================================================

\begin{schema}{BeginResponse}
  LuxFrame \\
  response? : RESPONSE
  \where
  sessionPhase = spProcessing \\
  toolPhase = tpNone \\
  activeSubagents = \emptyset \\
  sessionPhase' = spResponding \\
  pendingResponse' = \{ response? \} \\
  pendingPrompt' = \emptyset \\
  sessionId' = sessionId \\
  permMode' = permMode \\
  toolPhase' = tpNone \\
  currentTool' = \emptyset \\
  toolsThisTurn' = toolsThisTurn \\
  activeSubagents' = activeSubagents \\
  context' = context
\end{schema}

\begin{schema}{FinishResponse}
  LuxFrame \\
  hookResult? : HookDecision \\
  responseContext? : CONTEXTCHUNK
  \where
  sessionPhase = spResponding \\
  hookResult? = hdAllow \\
  sessionPhase' = spIdle \\
  pendingResponse' = \emptyset \\
  context' = (1 \upto maxContext) \dres (context \cat \langle responseContext? \rangle) \\
  toolsThisTurn' = 0 \\
  sessionId' = sessionId \\
  permMode' = permMode \\
  toolPhase' = tpNone \\
  currentTool' = \emptyset \\
  activeSubagents' = activeSubagents \\
  pendingPrompt' = \emptyset
\end{schema}

\begin{schema}{ForceContinue}
  LuxFrame \\
  hookResult? : HookDecision
  \where
  sessionPhase = spResponding \\
  hookResult? = hdBlock \\
  sessionPhase' = spProcessing \\
  pendingResponse' = \emptyset \\
  sessionId' = sessionId \\
  permMode' = permMode \\
  toolPhase' = tpNone \\
  currentTool' = \emptyset \\
  toolsThisTurn' = toolsThisTurn \\
  activeSubagents' = activeSubagents \\
  context' = context \\
  pendingPrompt' = pendingPrompt
\end{schema}

\begin{schema}{CompactContext}
  LuxFrame \\
  newLen? : \nat
  \where
  sessionPhase = spIdle \\
  \# context > 0 \\
  newLen? < \# context \\
  newLen? \leq maxContext \\
  context' = (1 \upto newLen?) \dres context \\
  sessionPhase' = sessionPhase \\
  sessionId' = sessionId \\
  permMode' = permMode \\
  toolPhase' = toolPhase \\
  currentTool' = currentTool \\
  toolsThisTurn' = toolsThisTurn \\
  activeSubagents' = activeSubagents \\
  pendingPrompt' = pendingPrompt \\
  pendingResponse' = pendingResponse
\end{schema}

\begin{schema}{EndSession}
  \Delta State \\
  reason? : SessionEndReason
  \where
  sessionPhase = spIdle \\
  activeSubagents = \emptyset \\
  sessionPhase' = spEnded \\
  sessionId' = \emptyset \\
  toolPhase' = tpNone \\
  currentTool' = \emptyset \\
  toolsThisTurn' = 0 \\
  activeSubagents' = \emptyset \\
  context' = \langle \rangle \\
  pendingPrompt' = \emptyset \\
  pendingResponse' = \emptyset \\
  permMode' = permMode \\
  %
  displayMode' = dmOff \\
  activeScene' = \emptyset \\
  sceneKind' = \emptyset \\
  pendingRecv' = zfalse \\
  beadsBoardStale' = zfalse
\end{schema}

% ============================================================
\section{System Invariants}
% ============================================================

\begin{enumerate}
  \item \textbf{Scene requires display}: \texttt{displayMode = dmOff
    $\implies$ activeScene = $\emptyset$}. No scene can exist when the
    display is disabled.

  \item \textbf{No orphan interactions}: \texttt{pendingRecv = ztrue
    $\implies$ activeScene $\neq \emptyset$}. Cannot block for
    interaction without an active scene.

  \item \textbf{Scene kind consistency}: \texttt{activeScene = $\emptyset$
    $\implies$ sceneKind = $\emptyset$} and
    \texttt{activeScene $\neq \emptyset$ $\implies$ $\#$sceneKind = 1}.

  \item \textbf{Beads staleness scoping}: \texttt{beadsBoardStale = ztrue
    $\implies$ sceneKind = \{skBeads\}}. Staleness only applies when the
    beads board is the active scene.

  \item \textbf{Display mode gates consumers}: \texttt{ShowScene} and
    \texttt{MarkBeadsBoardStale} require \texttt{displayMode = dmOn}.
    When lux is off, all visual operations are refused.

  \item \textbf{Clean session boundaries}: all lux state cleared on
    \texttt{EndSession}.
\end{enumerate}

% ============================================================
\section{Consumer Integration Patterns}
% ============================================================

Lux is a building block. These patterns describe how consumers
discover and use lux, following \texttt{integration.md}'s tiered
protocol.

\subsection{L0: Presence Detection}

Consumers check for lux via MCP tool availability (L1 discovery):

\begin{verbatim}
  # In consumer's hook handler:
  # Check if lux MCP tools are available in this session
  # If available AND displayMode = dmOn: render visual output
  # If not: degrade to text-only
\end{verbatim}

The \texttt{displayMode} flag is the L3 state that consumers read.
When \texttt{/lux y} is called, it writes to a config file that
consumers can check.

\subsection{Beads Integration}

\textbf{Trigger}: \texttt{PostToolUse} on \texttt{Bash} commands
matching \texttt{bd create|update|close}.

\textbf{Pattern}: When lux is enabled and the beads board is active,
mark it stale. Claude sees the staleness context and re-renders the
board with fresh data.

\textbf{Degradation}: Without lux, \texttt{bd ready} outputs text.
With lux, it renders a filterable board with search, filters, and
detail panel.

\subsection{Quarry Integration}

\textbf{Trigger}: \texttt{PostToolUse} on \texttt{quarry find}.

\textbf{Pattern}: When lux is enabled, render search results as an
interactive data explorer (filterable table with document detail panel).
Without lux, results are inline text.

\subsection{PR Workflow Integration}

\textbf{Trigger}: \texttt{PostToolUse} on
\texttt{mcp\_\_github\_\_create\_pull\_request}.

\textbf{Pattern}: When lux is enabled, display a PR dashboard
showing status, checks, review comments. Without lux, PR status is
text-only in the conversation.

\subsection{Z-Spec Integration}

\textbf{Trigger}: After \texttt{/z-spec:test} or \texttt{/z-spec:check}.

\textbf{Pattern}: When lux is enabled, display model-check results
as a dashboard (states explored, transitions fired, operations covered,
invariant status). The \texttt{claude-code.lux} file we created earlier
is an example of this pattern.

\subsection{Session Dashboard}

\textbf{Trigger}: \texttt{SessionStart} or explicit \texttt{/lux dashboard}.

\textbf{Pattern}: Compose a session overview from multiple sources:
git status (working tree state), beads (active issues), biff
(unread messages, active wall), vox (notify mode). Each data source
degrades independently.

% ============================================================
\section{Hook Event Map}
% ============================================================

\begin{tabular}{lllll}
\hline
\textbf{Hook} & \textbf{Matcher} & \textbf{Operation} & \textbf{Lux Effect} & \textbf{Blocks?} \\
\hline
SessionStart & * & StartSession & Load display config & No \\
PostToolUse & lux\_\_* & (suppress output) & Format panel & No \\
PostToolUse & Bash (bd *) & MarkBeadsBoardStale & Flag refresh & No \\
PostToolUse & quarry\_\_find & (consumer) & Data explorer & No \\
PostToolUse & github\_\_*\_pr & (consumer) & PR dashboard & No \\
Stop & * & (consumer) & Summary visual & No \\
SessionEnd & * & EndSession & Clear display & No \\
\hline
\end{tabular}

\textbf{Note}: Lux itself registers only SessionStart and PostToolUse
(suppress-output) hooks. The consumer integration hooks (Bash bd,
quarry find, github PR) are registered by the \emph{consuming} plugins
or by session-level configuration. Lux as a building block has no
upstream awareness.

% ============================================================
\section{Implementation Validation}
% ============================================================

\subsection{Built and Correct}

\begin{tabular}{llll}
\hline
\textbf{Component} & \textbf{Status} & \textbf{Notes} \\
\hline
SessionStart hook & Correct & Deploys commands, auto-allows tools \\
PostToolUse suppress & Correct & Two-channel display for all MCP tools \\
MCP tools (show, update, clear, recv, ping) & Correct & 10 tools exposed \\
Beads board skill & Correct & /lux:beads renders filterable board \\
Dashboard skill & Correct & /lux:dashboard composes metrics+charts \\
Data explorer skill & Correct & /lux:data-explorer with built-in filters \\
Diagram skill & Correct & /lux:diagram with auto-layout \\
\hline
\end{tabular}

\subsection{Not Yet Built}

\begin{tabular}{llll}
\hline
\textbf{Feature} & \textbf{Z Operation} & \textbf{Effort} & \textbf{Notes} \\
\hline
/lux y and /lux n toggle & EnableDisplay/DisableDisplay & Medium & Needs config file + command \\
Display mode config & (L3 state) & Low & .lux/config.md with displayMode \\
Beads board auto-refresh & MarkBeadsBoardStale & Medium & PostToolUse on Bash(bd *) \\
Quarry search display & (consumer pattern) & Low & PostToolUse on quarry find \\
PR dashboard & (consumer pattern) & Medium & PostToolUse on github PR \\
Z-spec results display & (consumer pattern) & Low & Already have .lux file pattern \\
Session dashboard & (consumer pattern) & Medium & Compose from multiple sources \\
\hline
\end{tabular}

\subsection{Design Note: Building Block Boundary}

Per \texttt{integration.md}, lux is a building block with no upstream
awareness. The consumer integration operations in this Z spec model
what \emph{consumers do} with lux, not what lux does itself. In
implementation:

\begin{itemize}
  \item \texttt{ShowScene}, \texttt{UpdateScene}, \texttt{ClearDisplay},
    \texttt{WaitForInteraction}, \texttt{ReceiveInteraction} are lux
    MCP tools --- implemented in lux.
  \item \texttt{EnableDisplay}, \texttt{DisableDisplay} are lux commands
    --- implemented in lux.
  \item \texttt{MarkBeadsBoardStale} would be a PostToolUse hook in the
    consuming plugin (biff or a shared integration hook), not in lux.
  \item Consumer display triggers (quarry search, PR dashboard) are
    implemented in the consumer's hooks or skills, not in lux.
\end{itemize}

This respects the one-way arrow: consumers discover lux, lux never
discovers consumers.

\end{document}
